VCD has three components, each designed to be useful independently but
engineered to compose. We describe them in the order they appear in a
single end-to-end invocation: schema verification, capability minting,
and gate enforcement (Figure 6).

\subsubsection{3.1 SMT-Backed Policy
Verification}\label{smt-backed-policy-verification}

Given a tool with JSON Schema \(S\) and a security policy \(P\), we ask:
does every JSON object satisfying \(S\) also satisfy \(P\)? If yes, the
policy is \emph{complete} over the schema; if no, there exists a
concrete counterexample tool call.

We support a fragment of JSON Schema that has decidable semantics under
the combined theory of strings, integers, and booleans in Z3
{[}@demoura2008z3{]}: top-level objects with primitive properties
(\texttt{string}, \texttt{number}, \texttt{integer}, \texttt{boolean}),
the \texttt{enum} keyword on string-valued properties,
\texttt{minimum}/\texttt{maximum} constraints on numeric properties, and
the \texttt{required} array. Constructs outside this fragment ---
recursive references, \texttt{oneOf}/\texttt{anyOf} schemas, arbitrary
regex constraints on string fields, unbounded arrays --- cause the
prover to raise \texttt{UnsupportedGrammar} rather than silently
approximate. Section 8 sketches the extensions needed to widen the
fragment.

We support a corresponding fragment of policy language. A policy is a
mapping from constraint kinds to per-field values:
\texttt{forbidden\_values:\ \{field\ →\ {[}v\_1,\ ...,\ v\_n{]}\}} and
\texttt{allowed\_values:\ \{field\ →\ {[}v\_1,\ ...,\ v\_n{]}\}} for
set-valued constraints on string fields; \texttt{max\_value} and
\texttt{min\_value} for numeric bounds;
\texttt{required\_present:\ {[}field,\ ...{]}} for additional presence
requirements beyond the schema's \texttt{required}; and
\texttt{forbidden\_field\_combinations:\ {[}{[}f\_1,\ f\_2{]},\ ...{]}}
for forbidden co-occurrences.

The encoding into Z3 is direct. Each schema property becomes an SMT
variable of matching sort. Each property's presence is a Boolean.
Schema-level constraints (enum membership, numeric bounds, required
presence) are conjuncted as preconditions. The policy is encoded as a
\emph{violation} disjunction: a single Z3 query asks whether there
exists an assignment satisfying the schema preconditions and at least
one violation. If \texttt{unsat}, the policy is proven complete. If
\texttt{sat}, Z3 produces a model from which we extract a concrete
counterexample call.

The prover emits a \texttt{ProofResult} containing canonical-JSON hashes
of the schema and policy, the proof status (\texttt{PROVEN},
\texttt{REFUTED}, or \texttt{UNDECIDED}), and any counterexample. The
proof result is itself hashed to produce a content-addressed
\texttt{proof\_hash} that capability tokens can cite.

\subsubsection{3.2 Capability Tokens}\label{capability-tokens}

A capability token \(t\) is an Ed25519-signed structure binding
\texttt{(agent\_id,\ tool,\ constraints,\ issued\_at,\ not\_before,\ expires\_at,\ parent\_id,\ policy\_proof\_hash,\ issuer,\ key\_id)}.
The signature covers a canonical-JSON serialisation of the bound fields.
The token's identifier is the SHA-256 of the same canonical body,
prefixed with \texttt{cap:}; the identifier is therefore
content-addressed and unforgeable absent collision.

The constraint schema mirrors the policy schema accepted by the prover,
by deliberate design: a token can cite the proof's hash and the gate can
statically verify that the token's constraints are at least as tight as
the proven policy.

\textbf{Attenuation.} Given a parent token \(p\) held by an authorised
principal, the issuer's \texttt{attenuate} primitive derives a child
token \(c\) with: (a) the same \texttt{tool} as \(p\); (b) an
\texttt{agent\_id} equal to or a prefix-extension of \(p\)'s; (c)
constraints that are \emph{narrower} along every dimension; (d) an
\texttt{expires\_at} no later than \(p\)'s; (e) \texttt{parent\_id} set
to \(p\)'s identifier. Narrowing is defined per constraint kind:

\begin{longtable}[]{@{}
  >{\raggedright\arraybackslash}p{(\linewidth - 2\tabcolsep) * \real{0.5000}}
  >{\raggedright\arraybackslash}p{(\linewidth - 2\tabcolsep) * \real{0.5000}}@{}}
\toprule\noalign{}
\begin{minipage}[b]{\linewidth}\raggedright
Kind
\end{minipage} & \begin{minipage}[b]{\linewidth}\raggedright
Narrowing operation
\end{minipage} \\
\midrule\noalign{}
\endhead
\bottomrule\noalign{}
\endlastfoot
\texttt{forbidden\_values{[}f{]}} & union:
\(c.\text{forbidden}_f \supseteq p.\text{forbidden}_f\) \\
\texttt{allowed\_values{[}f{]}} & intersection:
\(c.\text{allowed}_f \subseteq p.\text{allowed}_f\) \\
\texttt{max\_value{[}f{]}} & minimum:
\(c.\text{max}_f \le p.\text{max}_f\) \\
\texttt{min\_value{[}f{]}} & maximum:
\(c.\text{min}_f \ge p.\text{min}_f\) \\
\texttt{required\_present} & superset:
\(c.\text{required} \supseteq p.\text{required}\) \\
\texttt{forbidden\_field\_combinations} & superset:
\(c.\text{forbidden\_combos} \supseteq p.\text{forbidden\_combos}\) \\
\end{longtable}

Asking \texttt{attenuate} for a constraint that would broaden
permissions silently retains the parent's tighter bound; asking for a
child lifetime that would outlive the parent raises an exception. Both
behaviours are mechanically enforced (Section 4).

The attenuation discipline is borrowed directly from the
object-capability literature {[}@miller2006robust{]} and from Macaroons
{[}@birgisson2014macaroons{]}, from which it differs primarily in the
cryptographic substrate (Ed25519 signatures rather than chained HMACs)
and the content-addressed identifier (preventing token-ID reassignment
attacks that bedevil bearer-token systems).

\subsubsection{3.3 The Gate}\label{the-gate}

The gate is the only path from agent intent to tool execution. It runs
out-of-process from both the model and the tool implementation, accepts
a \texttt{(token,\ tool,\ args)} triple, and returns a
\texttt{GateDecision} of \texttt{ALLOW} or \texttt{DENY} with a
structured reason. The gate performs eight checks in order, failing
closed on any one of them:

\begin{enumerate}
\def\labelenumi{\arabic{enumi}.}
\tightlist
\item
  \textbf{Issuer pinning.} The token's \texttt{key\_id} must appear in
  the gate's trusted-issuer map; otherwise DENY.
\item
  \textbf{Signature verification.} Verify the Ed25519 signature against
  the canonical body using the pinned public key.
\item
  \textbf{Identifier binding.} Recompute \texttt{cap:} + SHA-256 of the
  canonical body and compare to the token's \texttt{token\_id}; reject
  on mismatch. This blocks an attack where an adversary substitutes a
  valid signature against a different body.
\item
  \textbf{Time bounds.} Reject if \texttt{now\ \textless{}\ not\_before}
  or \texttt{now\ ≥\ expires\_at}.
\item
  \textbf{Tool match.} Reject if
  \texttt{token.tool\ ≠\ requested\_tool}.
\item
  \textbf{Agent scope.} Reject if the caller's claimed
  \texttt{agent\_id} is neither equal to the token's \texttt{agent\_id}
  nor a prefix-extension of it.
\item
  \textbf{Constraint satisfaction.} For each constraint kind in the
  token, evaluate the constraint against the supplied \texttt{args}.
  Reject on any violation.
\item
  \textbf{Chain resolution (optional).} If the token carries a
  \texttt{parent\_id}, walk the parent chain. Every intermediate token
  must verify under the gate's trusted-issuer map. Reject on any
  unresolved or invalid parent.
\end{enumerate}

The gate is a deterministic, side-effect-free decision procedure. Every
invocation emits a structured event suitable for consumption by an
append-only audit log; the gate does not itself maintain state.

\textbf{Gate-monopoly is an architectural assumption, not a mathematical
theorem.} Theorem 2 (§4) establishes that any call returning ALLOW from
the gate satisfies the in-force constraints; it does \emph{not}
establish that every tool invocation must go through the gate. That
latter property is the deployer's responsibility and is enforced
architecturally: tool implementations live in a separate process from
the agent runtime, exposed only behind a gate-mediated RPC interface,
with no eval/exec/direct-import paths reachable from agent-controlled
data. The reference implementation ships the gate as an out-of-process
server behind a Unix socket or local HTTPS endpoint, with the issuer key
in an HSM. A deployment that violates this discipline --- for example,
by linking the agent runtime and the tool implementation into one
address space and letting the agent reach the tool symbol directly ---
loses the property regardless of the gate's correctness. We call this
property out explicitly because §4's mechanisation provides no formal
guarantee about it.

\subsubsection{3.4 Composition}\label{composition}

The three primitives compose as follows. A platform operator authors a
tool schema and a policy, runs the SMT prover, and on \texttt{PROVEN}
obtains a \texttt{proof\_hash}. The operator mints a root capability
token citing the proof hash, signing with their issuer key. Tokens are
distributed to authorised agent contexts, possibly via attenuation
through orchestration layers. At tool-call time, the agent presents its
current token to the gate; the gate enforces. Every gate decision is
recorded into a hash-chained audit log signed by a separate operational
key, providing tamper-evident attribution that is independently
verifiable offline.

The four hashes --- schema, policy, proof, audit-leaf --- and three
signatures --- issuer (on the token), audit (on the chain checkpoint),
gate-decision (optional, on the structured event) --- close the trust
graph. Every link is content-addressed; every link is verifiable by any
party in possession of the corresponding public key. Figure 1 sketches
the resulting graph.

\begin{verbatim}
                  ┌─────────────────┐
                  │  JSON Schema S  │── sha256 ──► grammar_hash
                  └─────────────────┘                   │
                                                        ▼
                  ┌─────────────────┐           ┌──────────────┐
                  │   Policy P      │── sha256 ─►│ proof_hash   │ (PROVEN/REFUTED)
                  └─────────────────┘           └──────────────┘
                                                        │
                                            cited by    │
                                                        ▼
            ┌──────────────────────────────────────────────────┐
            │  Capability Token   (Ed25519, content-addressed) │
            │  (agent, tool, constraints, nbf, exp,            │
            │   policy_proof_hash, parent_id)                  │
            └──────────────────────────────────────────────────┘
                                │                       │
                       attenuate│                       │ presented at call site
                                ▼                       ▼
                       ┌────────────────┐    ┌────────────────────┐
                       │  Child Token   │    │  Gate.check(...)   │
                       │  parent_id=…   │    │  ALLOW or DENY     │
                       └────────────────┘    └────────────────────┘
                                                        │
                                              audit event
                                                        ▼
                                            ┌──────────────────────┐
                                            │  Hash-chained log    │
                                            │  Merkle root, signed │
                                            └──────────────────────┘

Figure 1: The trust graph. Every solid arrow is a content-addressed
reference; every signed artefact is verifiable offline against a pinned
issuer or audit key.
\end{verbatim}
